\documentclass[a4paper,11pt,oneside]{article}
% Modifica le impostazioni qui e compila una prova alla volta.
\usepackage[T1]{fontenc}
\usepackage[utf8]{inputenc}
\usepackage[italian]{babel}
\usepackage{lmodern}
\usepackage{amsmath,amssymb}
\usepackage[a4paper,margin=2.5cm,headheight=15pt]{geometry}
\usepackage{xcolor,graphicx}
\usepackage{tikz}
\usetikzlibrary{arrows.meta,calc}
\usepackage{booktabs,tabularx,array}
\usepackage{enumitem}
\usepackage{fancyhdr,multicol,listings}
\usepackage{makeidx}
\usepackage[hidelinks]{hyperref}
\makeindex

\definecolor{verde}{HTML}{176756}
\definecolor{grigiochiaro}{HTML}{F2F5F4}
\definecolor{blu}{HTML}{24598A}
\pagestyle{fancy}
\fancyhf{}
\fancyhead[L]{\small Laboratorio \LaTeX}
\fancyhead[R]{\small 2BS}
\fancyfoot[C]{\thepage}
\setlength{\parindent}{14pt}
\setlength{\parskip}{3pt}
\setcounter{tocdepth}{1}
\lstset{
 language=[LaTeX]TeX,
 basicstyle=\small\ttfamily,
 keywordstyle=\color{blu},
 commentstyle=\color{verde},
 backgroundcolor=\color{grigiochiaro},
 frame=single,rulecolor=\color{gray},
 breaklines=true,columns=fullflexible,
 showstringspaces=false,keepspaces=true,
 aboveskip=8pt,belowskip=8pt
}
% Contatore usato nella sezione dedicata.
\newcounter{esempio}
\newcommand{\esempio}{\refstepcounter{esempio}\par
  \textbf{Esempio \theesempio}\quad}
% Questi due comandi servono solo a rendere riconoscibili le attività.
\newcommand{\prova}{\subsection*{Riproduci e modifica}}
\newcommand{\controlla}{\paragraph{Controllo nel PDF}}
\title{Scrivere e organizzare con \LaTeX\\
\large Laboratorio guidato sulla classe article}
\author{Nome e cognome}
\date{\today}

\begin{document}
\maketitle
\thispagestyle{plain}
\section*{Come usare questo laboratorio}
Questo file contiene esempi già funzionanti. Leggi una sezione,
ritrova nel sorgente i comandi che producono il risultato e svolgi
le richieste imitando il modello. Inserisci le tue prove subito dopo
la consegna, lasciando visibile l'esempio. Compila dopo ogni modifica.

Il percorso è diviso in argomenti brevi. Puoi riprenderlo in più lezioni.
Per iniziare lavora su titolo, struttura, paragrafi, enfasi ed elenchi.
Poi passa a indice, riferimenti e impostazioni della pagina.
Tabelle, figure e strumenti avanzati completano il percorso.

\begin{enumerate}
 \item Sostituisci il nome dell'autore nel preambolo.
 \item Salva una copia personale del progetto.
 \item Per ogni argomento completa almeno la prima richiesta.
 \item Usa il controllo finale per verificare il risultato.
 \item Consegna il sorgente e il PDF aggiornati.
\end{enumerate}
Il sorgente contiene già i pacchetti necessari per gli esempi attivi.
La sezione su \texttt{minted} contiene una prova facoltativa da attivare
in un progetto separato. La prova sulle immagini permette di caricare
un proprio file e mostra un riquadro finché il file manca.

\tableofcontents
\section{La classe article e il preambolo}\label{sec:article}
\index{article}
\subsection{Le parti del file}
La prima riga sceglie la classe del documento e le sue opzioni.
Il preambolo termina quando inizia il corpo del documento.
I pacchetti aggiungono funzioni, come tabelle più comode, colori o disegni.

\begin{lstlisting}
\documentclass[a4paper,11pt,oneside]{article}
\usepackage[italian]{babel}
\begin{document}
Il testo del documento.
\end{document}
\end{lstlisting}
Questo blocco mostra un file minimo. Non incollarlo intero dentro il corpo
del progetto che stai usando, perché il tuo documento ha già un inizio
e una fine.

\subsection{Commenti e impostazioni}
Nel sorgente una riga che inizia con \texttt{\%} è un commento.
L'opzione \texttt{11pt} stabilisce la dimensione base dei caratteri.
Il formato A4 e i margini sono impostati anche nel pacchetto
\texttt{geometry}, che userai più avanti.

\prova
\begin{enumerate}
 \item Aggiungi nel preambolo un commento con il tuo nome e la classe.
 \item Cambia \texttt{11pt} in \texttt{12pt}, compila e confronta una pagina.
 \item Ripristina \texttt{11pt} e spiega in un commento che cosa è cambiato.
\end{enumerate}
\controlla Il commento non compare. Il cambiamento della dimensione base
può modificare anche il numero delle righe e delle pagine.

\clearpage
\section{Titolo, autore e data}\label{sec:titolo}
\subsection{Informazioni nel preambolo}
Titolo, autore e data sono dichiarati prima del corpo.
Il comando che stampa queste informazioni si trova all'inizio del testo.
\begin{lstlisting}
\title{Una relazione sui triangoli}
\author{Nome Cognome}
\date{\today}
% Nel corpo del documento:
\maketitle
\end{lstlisting}
Il titolo di questo laboratorio è costruito proprio così.
La data può essere automatica, scritta per esteso oppure lasciata vuota
con \verb|\date{}|.

\subsection{Un sottotitolo}
Un cambio di riga e un comando di dimensione permettono di aggiungere
un sottotitolo. Nel preambolo del tuo file trovi già questo esempio.
\begin{lstlisting}
\title{Una relazione sui triangoli\\
\large Esempi e costruzioni}
\end{lstlisting}

\prova
\begin{enumerate}
 \item Personalizza il titolo, il sottotitolo e l'autore del file.
 \item Prova prima la data automatica e poi una data scritta da te.
 \item Verifica l'effetto di una data vuota e scegli la versione finale.
\end{enumerate}
\controlla Le modifiche devono comparire nella pagina iniziale.
Un solo comando per stampare il titolo è sufficiente.

\clearpage
\section{Sezioni, sottosezioni e titoli senza numero}\label{sec:gerarchia}
\subsection{La gerarchia}
Questa è una sottosezione della sezione sui titoli.
\subsubsection{Un livello ulteriore}
Una sottosottosezione approfondisce un punto.
\paragraph{Un titolo nel paragrafo}
Questo livello introduce una porzione breve di testo.

\begin{lstlisting}
\section{Triangoli}
\subsection{Classificazione}
\subsubsection{In base ai lati}
\paragraph{Il triangolo isoscele}
Un triangolo con due lati congruenti.
\end{lstlisting}

\subsection{La versione con asterisco}
Il titolo seguente usa un asterisco.
\subsection*{Una precisazione}
Questo titolo non ha numero e non entra automaticamente nell'indice.
\begin{lstlisting}
\subsection*{Una precisazione}
\end{lstlisting}
In \texttt{article} i titoli principali sono le sezioni.
Il comando \verb|\chapter| appartiene a classi come \texttt{report}
e \texttt{book}.

\prova
\begin{enumerate}
 \item Aggiungi una sottosezione su un argomento matematico.
 \item Al suo interno inserisci due sottosottosezioni e un paragrafo titolato.
 \item Aggiungi un titolo con asterisco e confrontane la numerazione.
\end{enumerate}
\controlla I numeri devono essere generati dai comandi.
Non inserirli nel testo del titolo.

\clearpage
\section{Paragrafi, righe e spazi nel sorgente}\label{sec:paragrafi}
\subsection{Due paragrafi}
Questo primo paragrafo introduce una definizione.
Nel sorgente continua sulla riga seguente, ma rimane lo stesso paragrafo.

Questo secondo paragrafo contiene un esempio.
La riga vuota nel sorgente separa le due idee.
\begin{lstlisting}
Questo e il primo paragrafo.
Continua sulla riga seguente del sorgente.

Questo e il secondo paragrafo.
\end{lstlisting}
La riga vuota può essere sostituita da \verb|\par|.
Molti spazi consecutivi nel sorgente diventano normalmente uno spazio.

\subsection{Un'interruzione di riga}
\begin{flushleft}
Istituto David Maria Turoldo\\
Classe 2BS\\
Laboratorio di informatica
\end{flushleft}
Nel blocco precedente le righe sono separate da \verb|\\|.
Per iniziare un nuovo paragrafo usa una riga vuota.

\prova
\begin{enumerate}
 \item Scrivi una definizione in un paragrafo e un esempio nel successivo.
 \item Duplica l'intestazione e sostituiscila con i dati della tua relazione.
 \item Aggiungi dieci spazi tra due parole nel sorgente e osserva il PDF.
\end{enumerate}
\controlla Verifica la differenza tra un nuovo paragrafo e una nuova riga.

\clearpage
\section{Grassetto, corsivo e sottolineatura}\label{sec:enfasi}
\subsection{Esempi nel testo}
La \textbf{condizione di esistenza} va controllata.
Un \textit{fattore comune} può essere raccolto.
Una parola \underline{importante} può essere sottolineata.
Un termine può essere in \textbf{\textit{grassetto e corsivo}}.
Un nome di comando può essere scritto in \texttt{carattere monospaziato}.
\begin{lstlisting}
\textbf{condizione di esistenza}
\textit{fattore comune}
\underline{importante}
\textbf{\textit{grassetto e corsivo}}
\texttt{carattere monospaziato}
\end{lstlisting}

\subsection{Una modifica locale}
{\large Questa frase è più grande.}
Questa frase torna alla dimensione normale.
Le graffe delimitano l'effetto del comando di dimensione.

\prova
\begin{enumerate}
 \item Scrivi una definizione con il termine definito in grassetto.
 \item Metti una precisazione in corsivo e sottolinea una sola parola.
 \item Aggiungi una frase con un termine in grassetto e corsivo.
 \item Prova \verb|\small|, \verb|\large| e \verb|\Large| dentro gruppi di graffe.
\end{enumerate}
\controlla Il formato delle parole vicine deve restare quello previsto.
Se tutta la pagina cambia dimensione, controlla i confini del gruppo.

\clearpage
\section{Elenchi puntati, numerati e descrittivi}\label{sec:elenchi}
\subsection{Tre modelli}
\begin{itemize}
 \item Controllare i denominatori.
 \item Verificare i segni.
\end{itemize}
\begin{enumerate}
 \item Scomporre i polinomi.
 \item Scrivere le condizioni.
\end{enumerate}
\begin{description}
 \item[Monomio] Prodotto tra un coefficiente e potenze di lettere.
 \item[Polinomio] Somma di monomi.
\end{description}
\begin{lstlisting}
\begin{itemize}
  \item Prima voce.
  \item Seconda voce.
\end{itemize}
% Sostituisci itemize con enumerate per numerare.
% In description scrivi una voce come:
% \item[Termine] Spiegazione.
\end{lstlisting}
Un elenco può contenere un altro elenco dentro una delle sue voci.

\prova
\begin{enumerate}
 \item Scrivi una procedura numerata per risolvere un'equazione frazionaria.
 \item Dentro l'ultima voce aggiungi due controlli in un elenco puntato.
 \item Crea un elenco descrittivo con tre termini di geometria.
 \item Inserisci una nuova azione al secondo posto della procedura.
\end{enumerate}
\controlla I numeri si aggiornano. Le voci dell'elenco interno devono
rientrare rispetto a quelle esterne.

\clearpage
\section{Indice generale e titoli nell'indice}\label{sec:indice}
\index{indice!generale}
L'indice iniziale nasce da \verb|\tableofcontents|.
Raccoglie i titoli del documento e i numeri di pagina.
Per aggiornarlo può servire una seconda compilazione.
La profondità dell'indice è regolata da \texttt{tocdepth}.
In questo file vale 1 e include solo le sezioni.
\begin{lstlisting}
% Nel corpo, nel punto in cui vuoi l'indice:
\tableofcontents
% Nel preambolo, per includere le sottosezioni:
\setcounter{tocdepth}{2}
\end{lstlisting}

\subsection*{Nota sull'indice}
\phantomsection
\addcontentsline{toc}{subsection}{Nota sull'indice}
Questo titolo senza numero è stato aggiunto esplicitamente all'indice.
Il collegamento viene preparato da \verb|\phantomsection|.
Il pacchetto \texttt{hyperref} gestisce i richiami cliccabili.
\begin{lstlisting}
\subsection*{Nota sull'indice}
\phantomsection
\addcontentsline{toc}{subsection}{Nota sull'indice}
\end{lstlisting}
\prova
\begin{enumerate}
 \item Includi nell'indice le sottosezioni cambiando la profondità.
 \item Aggiungi una sottosezione e verifica che compaia.
 \item Crea un titolo senza numero e inseriscilo nell'indice con il modello.
\end{enumerate}
\controlla Nome del titolo e numero di pagina devono corrispondere.
Tieni presente che l'indice più lungo sposta le pagine successive.

\clearpage
\section{Riferimenti a sezioni ed equazioni}\label{sec:rimandi}
\subsection{Etichette e richiami}
Gli elenchi della sezione \ref{sec:elenchi} si trovano a pagina
\pageref{sec:elenchi}. I numeri derivano dall'etichetta della sezione.
\begin{lstlisting}
\section{Elenchi}
\label{sec:elenchi}
% Nel testo:
Vedi la sezione \ref{sec:elenchi}
a pagina \pageref{sec:elenchi}.
\end{lstlisting}

\subsection{Un'equazione numerata}
\begin{equation}
(a-b)^2=a^2-2ab+b^2
\label{eq:binomio}
\end{equation}
La formula \eqref{eq:binomio} è il quadrato di una differenza.
\begin{lstlisting}
\begin{equation}
(a-b)^2=a^2-2ab+b^2
\label{eq:binomio}
\end{equation}
La formula \eqref{eq:binomio}...
\end{lstlisting}
\prova
\begin{enumerate}
 \item Crea un'equazione numerata con una scomposizione.
 \item Assegna un'etichetta nuova e richiamala in una frase.
 \item Richiama anche una delle tue sezioni con il suo numero e la sua pagina.
\end{enumerate}
\controlla Ogni etichetta deve avere un nome unico.
Un'etichetta incollata insieme a un esempio va rinominata anche nei richiami.

\clearpage
\section{Rientri e distanza tra paragrafi}\label{sec:rientri}
La lunghezza \verb|\parindent| regola il rientro della prima riga.
La lunghezza \verb|\parskip| regola lo spazio tra paragrafi.
Un gruppo limita una modifica a una zona del testo.

\begingroup
\setlength{\parindent}{0pt}
\setlength{\parskip}{7pt}
Questo paragrafo è scritto con rientro nullo e distanza aumentata.

Questo secondo paragrafo permette di osservare la distanza.
Il testo riempie una parte della riga per rendere visibile il confronto.
\par
\endgroup

Questo testo è di nuovo fuori dal gruppo.
\begin{lstlisting}
\begingroup
\setlength{\parindent}{0pt}
\setlength{\parskip}{7pt}
Primo paragrafo.

Secondo paragrafo.
\par
\endgroup
\end{lstlisting}

Il primo paragrafo dopo un titolo, nelle impostazioni usuali di
\texttt{article}, comincia senza rientro.
Il comando \verb|\noindent| elimina il rientro del paragrafo che segue.

\prova
\begin{enumerate}
 \item Duplica il gruppo e usa un rientro di 8 mm e uno spazio di 2 pt.
 \item Scrivi tre paragrafi per rendere visibile l'effetto.
 \item Fai iniziare il terzo senza rientro usando \verb|\noindent|.
\end{enumerate}
\controlla Controlla il secondo paragrafo per valutare il rientro.
Fuori dal gruppo le impostazioni originali devono tornare attive.

\clearpage
\section{Allineare il testo}\label{sec:allineamento}
\begin{flushleft}
Un testo allineato a sinistra può avere il bordo destro irregolare.
Questo periodo è abbastanza lungo da occupare più di una riga.
\end{flushleft}
\begin{center}
Un titolo o un breve messaggio\\
può essere centrato.
\end{center}
\begin{flushright}
Un nome e una firma\\
possono essere allineati a destra.
\end{flushright}

\begin{lstlisting}
\begin{flushleft}
Testo a sinistra.
\end{flushleft}
\begin{center}
Testo centrato.
\end{center}
\begin{flushright}
Testo a destra.
\end{flushright}
\end{lstlisting}
Il comando \verb|\centering| si usa spesso dentro una figura o
una tabella. La sua azione va contenuta nel relativo ambiente o gruppo.

\prova
\begin{enumerate}
 \item Prepara una breve convocazione con destinatario a sinistra,
 titolo al centro e firma a destra.
 \item Scrivi un paragrafo lungo e confronta la giustificazione normale
 con \texttt{flushleft}.
\end{enumerate}
\controlla Il cambio di allineamento deve interessare il testo scelto.

\clearpage
\section{Lunghezze e spazi}\label{sec:lunghezze}
\index{lunghezze}
Una lunghezza può essere assoluta, come \texttt{1cm}, oppure riferita
allo spazio disponibile, come \verb|0.5\linewidth|.
\verb|\textwidth| è la larghezza del corpo principale della pagina.
\verb|\linewidth| è la larghezza disponibile nel contesto corrente.

\noindent Inizio\hspace{12mm}fine.

\vspace{4mm}
\noindent Questa frase segue uno spazio verticale.
\begin{lstlisting}
Inizio\hspace{12mm}fine.

\vspace{4mm}
Testo successivo.
\end{lstlisting}

\subsection{Due riquadri proporzionati}
\noindent\fbox{\parbox{0.35\linewidth}{Un riquadro stretto.}}
\hfill
\fbox{\parbox{0.45\linewidth}{Un riquadro più largo, con testo che
può andare a capo.}}

\subsection{Tenere due elementi insieme}
La scrittura \verb|pagina~12| usa uno spazio che impedisce di separare
le due parole a fine riga. Il carattere \verb|~| è utile anche tra
un'abbreviazione e il numero che segue.

\prova
\begin{enumerate}
 \item Porta lo spazio orizzontale a 2 cm e quello verticale a 8 mm.
 \item Duplica i riquadri e cambia le proporzioni mantenendoli sulla stessa riga.
 \item Aggiungi una frase che usi uno spazio non separabile.
\end{enumerate}
\controlla La somma delle larghezze, degli spazi e dei bordi deve
rientrare nella riga. Controlla eventuali avvisi \texttt{Overfull}.

\clearpage
\section{Formato del foglio e margini}\label{sec:margini}
La classe usa A4. Il pacchetto \texttt{geometry} imposta margini e
spazio per l'intestazione. Le misure modificano lo spazio utile
e quindi il modo in cui il testo va a capo.
\begin{lstlisting}
% Questa riga sostituisce quella di geometry nel preambolo.
\usepackage[
  a4paper,
  left=2.8cm,
  right=2.8cm,
  top=2.5cm,
  bottom=2.5cm,
  headheight=15pt
]{geometry}
\end{lstlisting}
Sostituisci la riga già presente quando provi il modello.
Caricare lo stesso pacchetto con opzioni incompatibili può causare un errore.

\subsection{Un controllo delle misure}
La larghezza corrente del testo è \the\textwidth.
L'altezza corrente del testo è \the\textheight.
Questi numeri sono stampati interrogando le lunghezze del documento.

\prova
\begin{enumerate}
 \item Annota il numero di pagine della tua versione.
 \item Imposta margini sinistro e destro di 3 cm.
 \item Confronta una pagina con una tabella e annota che cosa cambia.
 \item Scegli le misure finali controllando che nulla esca dai margini.
\end{enumerate}
\controlla Anche una tabella o un'immagine prima corretta può diventare
troppo larga dopo una modifica dei margini.

\clearpage
\section{Intestazioni e piè di pagina}\label{sec:intestazioni}
Questo documento usa \texttt{fancyhdr}. La sigla \texttt{L} indica
la sinistra, \texttt{C} il centro e \texttt{R} la destra.
\begin{lstlisting}
% Impostazioni nel preambolo
\pagestyle{fancy}
\fancyhf{}
\fancyhead[L]{Laboratorio}
\fancyhead[R]{2BS}
\fancyfoot[C]{\thepage}
\end{lstlisting}
\verb|\fancyhf{}| svuota i campi prima di impostarli.
Il comando \verb|\thepage| stampa il numero corrente.

\subsection{Lo stile di una singola pagina}
La pagina iniziale usa \verb|\thispagestyle{plain}|.
Puoi usare \verb|\thispagestyle{empty}| per togliere intestazione
e piè di pagina solo dalla pagina corrente.

\prova
\begin{enumerate}
 \item Sostituisci l'intestazione sinistra con il tuo cognome.
 \item Sposta il numero di pagina a destra.
 \item Inserisci la classe nel piè di pagina sinistro.
 \item Prova lo stile \texttt{empty} su questa pagina, osserva l'effetto
 e poi ripristina lo stile iniziale.
\end{enumerate}
\controlla La modifica di una singola pagina non deve cancellare
le intestazioni di tutte le pagine.

\clearpage
\section{Numerazione e cambi di pagina}\label{sec:pagine}
\subsection{Numeri arabi e romani}
Le pagine di questo file sono numerate con cifre arabe.
Il modello seguente si prova in un progetto separato perché cambia
la numerazione durante il documento.
\begin{lstlisting}
\documentclass{article}
\begin{document}
\pagenumbering{roman}
Introduzione.
\clearpage
\pagenumbering{arabic}
Prima pagina del testo principale.
\end{document}
\end{lstlisting}
\verb|\pagenumbering| cambia il formato e riporta il contatore a 1.
\verb|\setcounter{page}{5}| assegna invece un valore preciso.

\subsection{Nuova pagina e oggetti in attesa}
\verb|\newpage| termina la pagina corrente.
\verb|\clearpage| termina la pagina e fa uscire anche figure e tabelle
che aspettano di essere collocate.

\prova
\begin{enumerate}
 \item Esegui il modello in un progetto di prova e verifica i numeri.
 \item Torna a questo file e trova i comandi tra una sezione e l'altra.
 \item Inserisci un cambio pagina in una tua sezione e osserva lo spazio rimasto.
\end{enumerate}
\controlla Evita di riportare la numerazione a 1 più volte nella versione
da consegnare. I riferimenti devono restare chiari.

\clearpage
\section{Una facciata e due facciate}\label{sec:facciate}
La prima riga contiene l'opzione \texttt{oneside}.
Per una prova fronte-retro puoi usare \texttt{twoside}.
In quel caso il margine interno è quello rivolto alla rilegatura.
\begin{lstlisting}
\documentclass[a4paper,11pt,twoside]{article}
% Nel pacchetto geometry sostituisci le misure con:
% inner=3cm,outer=2cm,top=2.5cm,bottom=2.5cm
% Per intestazioni alternate:
\fancyhead[LE,RO]{2BS}
\fancyhead[LO,RE]{Relazione}
\end{lstlisting}
Le lettere \texttt{E} e \texttt{O} indicano pagine pari e dispari.
Questo esempio va provato dopo aver attivato \texttt{twoside}.

\prova
\begin{enumerate}
 \item Salva una copia del progetto con un nome diverso.
 \item Attiva \texttt{twoside} e imposta margini interno ed esterno diversi.
 \item Svuota i campi delle intestazioni e usa il modello alternato.
 \item Confronta due pagine consecutive a sinistra e a destra.
\end{enumerate}
\controlla Il margine interno cambia lato tra pagina pari e pagina dispari.
Osserva il documento come se fosse già rilegato.

\clearpage
\section{Costruire una tabella}\label{sec:tabelle}
\index{tabelle}
Una tabella ha colonne allineate e righe di dati.
In \texttt{tabular}, \texttt{l}, \texttt{c} e \texttt{r}
indicano rispettivamente sinistra, centro e destra.
\begin{center}
\begin{tabular}{lcr}
 \hline
 Figura & Lati & Perimetro\\
 \hline
 Equilatero & $a,a,a$ & $3a$\\
 Isoscele & $a,a,b$ & $2a+b$\\
 \hline
\end{tabular}
\end{center}
\begin{lstlisting}
\begin{tabular}{lcr}
\hline
Figura & Lati & Perimetro\\
\hline
Equilatero & $a,a,a$ & $3a$\\
Isoscele & $a,a,b$ & $2a+b$\\
\hline
\end{tabular}
\end{lstlisting}
Il simbolo \verb|&| separa le celle e \verb|\\| termina la riga.

\subsection{Colonne che accolgono testo}
\begin{tabularx}{\linewidth}{lX}
\toprule
Termine & Spiegazione\\
\midrule
Altezza & Segmento perpendicolare alla retta di un lato condotto
dal vertice opposto.\\
Mediana & Segmento che unisce un vertice con il punto medio del lato opposto.\\
\bottomrule
\end{tabularx}
La colonna \texttt{X} di \texttt{tabularx} occupa lo spazio rimasto
e permette al testo di andare a capo.

\prova
\begin{enumerate}
 \item Aggiungi alla prima tabella un triangolo scaleno.
 \item Duplica la seconda tabella e aggiungi la definizione di bisettrice.
 \item Prova una tabella con quattro colonne e almeno tre righe di dati.
\end{enumerate}
\controlla Ogni riga deve avere il numero di celle previsto.

\clearpage
\section{Didascalie, posizione e riferimenti delle tabelle}\label{sec:float}
L'ambiente \texttt{table} contiene la tabella e la sua didascalia.
Il suo numero può essere richiamato dal testo.
\begin{table}[htbp]
\centering
\caption{Un confronto tra due scomposizioni.}
\label{tab:scomposizioni}
\begin{tabular}{ll}
\toprule
Polinomio & Scomposizione\\
\midrule
$x^2-9$ & $(x-3)(x+3)$\\
$x^2+4x+4$ & $(x+2)^2$\\
\bottomrule
\end{tabular}
\end{table}
La tabella \ref{tab:scomposizioni} usa la didascalia sopra i dati.
\begin{lstlisting}
\begin{table}[htbp]
\centering
\caption{Un titolo breve.}
\label{tab:nuova}
% Inserisci qui il tuo ambiente tabular.
\end{table}
\end{lstlisting}
Le lettere suggeriscono collocazioni possibili.
\texttt{h} indica qui, \texttt{t} in alto, \texttt{b} in basso e
\texttt{p} una pagina dedicata agli oggetti.
Sono preferenze che il sistema valuta insieme allo spazio disponibile.

\prova
\begin{enumerate}
 \item Trasforma una tua tabella in un oggetto con didascalia ed etichetta.
 \item Richiamala in una frase con \verb|\ref|.
 \item Prova prima \texttt{[t]} e poi \texttt{[htbp]}.
 \item Osserva dove compare rispetto al paragrafo che la precede.
\end{enumerate}
\controlla L'etichetta deve essere dopo la didascalia.
Il numero nel richiamo deve coincidere con quello della tabella.

\clearpage
\section{Un disegno con TikZ}\label{sec:tikz}
\index{TikZ}
Con TikZ il disegno è descritto nel sorgente.
Le coordinate fissano punti; le istruzioni tracciano segmenti e nomi.
\begin{center}
\begin{tikzpicture}[scale=0.8]
 \coordinate (A) at (0,0);
 \coordinate (B) at (6,0);
 \coordinate (C) at (2,3);
 \coordinate (H) at (2,0);
 \draw[thick] (A)--(B)--(C)--cycle;
 \draw[dashed,verde] (C)--(H);
 \node[below left] at (A) {$A$};
 \node[below right] at (B) {$B$};
 \node[above] at (C) {$C$};
 \node[below] at (H) {$H$};
\end{tikzpicture}
\end{center}
\begin{lstlisting}
\begin{tikzpicture}[scale=0.8]
\coordinate (A) at (0,0);
\coordinate (B) at (6,0);
\coordinate (C) at (2,3);
\draw[thick] (A)--(B)--(C)--cycle;
\node[above] at (C) {$C$};
\end{tikzpicture}
\end{lstlisting}
Il comando \texttt{cycle} chiude il percorso.
Il punto e virgola conclude ogni istruzione TikZ.

\prova
\begin{enumerate}
 \item Duplica il disegno e sposta $C$ in $(3,3)$.
 \item Aggiorna anche $H$ e confronta i due triangoli.
 \item Cambia colore e tratteggio dell'altezza.
 \item Scrivi una frase che motivi perché il secondo triangolo è isoscele.
\end{enumerate}
\controlla L'altezza deve essere perpendicolare alla base.

\clearpage
\section{Figure, immagini e dimensioni}\label{sec:figure}
L'ambiente \texttt{figure} aggiunge didascalia e numero a un disegno.
L'esempio seguente funziona senza caricare file esterni.
\begin{figure}[htbp]
\centering
\begin{tikzpicture}
 \draw[thick,blu] (0,0) rectangle (4,2);
 \draw[dashed] (0,0)--(4,2);
 \node at (2,-0.4) {$b=4$};
 \node[right] at (4,1) {$h=2$};
\end{tikzpicture}
\caption{Un rettangolo diviso dalla diagonale.}
\label{fig:rettangolo}
\end{figure}
La figura \ref{fig:rettangolo} può essere richiamata come una tabella.

\subsection{Caricare una propria immagine}
Il blocco nel sorgente usa il file \texttt{immagine.png}, se presente.
Finché non lo carichi mostra un riquadro.
\begin{figure}[htbp]
\centering
\IfFileExists{immagine.png}{
 \includegraphics[width=0.45\linewidth]{immagine.png}
}{
 \fbox{\parbox[c][18mm][c]{0.55\linewidth}{\centering
 Carica una tua immagine chiamata \texttt{immagine.png}.}}
}
\caption{La mia immagine di prova.}
\label{fig:immagine}
\end{figure}
\prova
\begin{enumerate}
 \item Carica un'immagine con il nome indicato e verifica che appaia.
 \item Prova larghezze \texttt{0.3}, \texttt{0.5} e \texttt{0.7}.
 \item Cambia la didascalia e aggiungi un riferimento alla figura nel testo.
\end{enumerate}
\controlla Rispetta maiuscole ed estensione del nome del file.

\clearpage
\section{Elenco delle figure e delle tabelle}\label{sec:elenchi-oggetti}
Questi elenchi raccolgono le didascalie e i numeri di pagina.
Dopo una modifica possono richiedere una nuova compilazione.
\listoffigures
\listoftables
\begin{lstlisting}
\listoffigures
\listoftables
\end{lstlisting}

\subsection{Una didascalia breve per l'elenco}
Puoi mostrare un titolo breve nell'elenco e una spiegazione più estesa
sotto la figura usando un argomento opzionale.
\begin{lstlisting}
\caption[Triangolo isoscele]{
Triangolo isoscele con l'altezza relativa alla base.}
\end{lstlisting}

\prova
\begin{enumerate}
 \item Aggiungi una figura o una tabella alla tua versione.
 \item Imposta una didascalia estesa con un titolo breve per l'elenco.
 \item Compila di nuovo e confronta l'elenco con l'oggetto nel testo.
\end{enumerate}
\controlla Il titolo breve compare nell'elenco, la didascalia estesa
resta accanto alla figura o alla tabella.

\clearpage
\section{Scrivere in più colonne}\label{sec:colonne}
L'ambiente \texttt{multicols} divide una zona del documento.
Il numero tra graffe stabilisce quante colonne creare.
\begin{multicols}{2}
\textbf{Definizione.}
Una mediana unisce un vertice del triangolo con il punto medio
del lato opposto.
\columnbreak
\textbf{Osservazione.}
Un triangolo ha tre mediane. La posizione del punto medio è
decisiva per riconoscere il segmento.
\end{multicols}
\begin{lstlisting}
\begin{multicols}{2}
Testo nella prima colonna.
\columnbreak
Testo nella seconda colonna.
\end{multicols}
\end{lstlisting}
Senza cambio forzato, il testo viene distribuito dal sistema.
Nel preambolo puoi cambiare la distanza tra le colonne.
\begin{lstlisting}
\setlength{\columnsep}{8mm}
\end{lstlisting}

\prova
\begin{enumerate}
 \item Crea una scheda a due colonne con una definizione e un esempio.
 \item Togli \verb|\columnbreak| e aggiungi testo per vedere la distribuzione.
 \item Prova tre colonne con tre brevi definizioni.
\end{enumerate}
\controlla Le formule e le immagini devono stare nella colonna.
Dentro una colonna usa \verb|\linewidth| per misurare la larghezza utile.

\clearpage
\section{Colori nel testo e nei riquadri}\label{sec:colori}
Il pacchetto \texttt{xcolor} permette di definire colori e applicarli.
Nel preambolo sono già definiti \texttt{verde}, \texttt{blu}
e \texttt{grigiochiaro}.
\begin{center}
\textcolor{verde}{Condizione verificata}\qquad
\textcolor{blu}{Controllo da eseguire}
\end{center}
\begin{center}
\colorbox{grigiochiaro}{\strut Una precisazione.}
\end{center}
\begin{lstlisting}
% Nel preambolo:
\definecolor{verde}{HTML}{176756}
% Nel testo:
\textcolor{verde}{Condizione verificata}
\colorbox{grigiochiaro}{Una precisazione.}
\end{lstlisting}
Il modello \texttt{HTML} usa una scrittura esadecimale.
Il modello \texttt{RGB} usa tre interi tra 0 e 255.
\begin{lstlisting}
\definecolor{arancio}{RGB}{190,95,15}
\end{lstlisting}

\prova
\begin{enumerate}
 \item Definisci un nuovo colore e applicalo a un termine.
 \item Inserisci un breve riquadro con un richiamo sulle condizioni di esistenza.
 \item Prova una stampa in scala di grigi o osserva il PDF senza colore.
\end{enumerate}
\controlla Il testo deve restare leggibile e le informazioni importanti
devono essere riconoscibili anche dalle parole usate.

\clearpage
\section{Note a piè di pagina e a margine}\label{sec:note}
\subsection{Un approfondimento a fondo pagina}
Una condizione di esistenza\footnote{Una frazione algebrica è definita
per i valori che rendono il denominatore diverso da zero.}
può essere spiegata in una nota.
\begin{lstlisting}
Una condizione di esistenza\footnote{
Spiegazione della condizione.}
\end{lstlisting}
La nota viene numerata automaticamente.
Inserire il comando in un normale paragrafo permette di osservarne
il collegamento con il punto preciso del testo.

\subsection{Un'annotazione laterale}
Questa riga ha una nota a margine.
\marginpar{\footnotesize Controlla il denominatore.}
Il comando colloca una precisazione nel margine.
Un testo troppo lungo può occupare molte righe.
\begin{lstlisting}
Testo principale.
\marginpar{\footnotesize Una nota breve.}
\end{lstlisting}

\prova
\begin{enumerate}
 \item Aggiungi una nota a piè di pagina a una definizione.
 \item Inserisci una nota laterale di tre o quattro parole.
 \item Aggiungi una nota prima di quella esistente e osserva i numeri.
\end{enumerate}
\controlla La nota laterale deve restare dentro il foglio.
Controlla il risultato anche dopo una modifica dei margini.

\clearpage
\section{Contatori e numerazione automatica}\label{sec:contatori}
\index{contatori}
Pagine, sezioni ed equazioni hanno contatori.
Puoi anche crearne uno per una tua sequenza.
Nel preambolo è definito il contatore \texttt{esempio}.
\esempio Il quadrato di $x+1$ è $x^2+2x+1$.\par
\esempio La differenza $x^2-1$ si scompone in $(x-1)(x+1)$.\par

\begin{lstlisting}
% Nel preambolo:
\newcounter{esempio}
\newcommand{\esempio}{\refstepcounter{esempio}\par
\textbf{Esempio \theesempio}\quad}
% Nel corpo:
\esempio Un primo testo.\par
\esempio Un secondo testo.\par
\end{lstlisting}
La definizione è già presente nel tuo preambolo.
\verb|\refstepcounter| incrementa e rende disponibile il valore
per un'etichetta.
\verb|\setcounter| assegna un valore;
\verb|\addtocounter| aggiunge una quantità.

\prova
\begin{enumerate}
 \item Aggiungi tre esempi con il comando predisposto.
 \item Assegna un'etichetta a un esempio e richiamalo con \verb|\ref|.
 \item Nella copia di prova imposta il contatore a 9 prima di un nuovo esempio.
\end{enumerate}
\controlla Il successivo esempio deve essere il numero 10.
Il richiamo deve usare il numero dell'esempio, non quello della sezione.

\clearpage
\section{Mostrare codice con verbatim e listings}\label{sec:codice}
Per stampare un breve comando puoi usare \verb|\verb|.
Per un blocco intero puoi usare \texttt{verbatim} oppure
\texttt{lstlisting}. Gli esempi di questo laboratorio usano il secondo.
\begin{verbatim}
\textbf{Qui vedo il comando completo.}
\end{verbatim}

\begin{lstlisting}[numbers=left]
% Questo commento compare nel PDF.
\begin{enumerate}
  \item Prima voce.
  \item Seconda voce.
\end{enumerate}
\end{lstlisting}
Il contenuto viene mostrato come codice.
Le impostazioni di \texttt{listings} nel preambolo scelgono caratteri,
bordi e colori.
\begin{lstlisting}
\lstset{
  language=[LaTeX]TeX,
  basicstyle=\small\ttfamily,
  keywordstyle=\color{blu},
  frame=single,
  breaklines=true
}
\end{lstlisting}

\prova
\begin{enumerate}
 \item Mostra il codice di una tua tabella in un nuovo blocco.
 \item Aggiungi \texttt{numbers=left} al blocco per numerarne le righe.
 \item Cambia colore delle parole chiave e dimensione del carattere.
\end{enumerate}
\controlla Nel blocco deve comparire il codice.
La tabella effettiva può stare sotto per confrontare sorgente e risultato.

\clearpage
\section{Approfondimento sul pacchetto minted}\label{sec:minted}
Il pacchetto \texttt{minted} evidenzia il codice usando un elaboratore
esterno. La disponibilità dipende dall'ambiente di compilazione.
Questa prova si esegue in un progetto separato.
Il file principale continua a usare \texttt{listings}.

\subsection{Il modello da provare}
\begin{lstlisting}
\documentclass{article}
\usepackage[T1]{fontenc}
\usepackage{minted}
\begin{document}
\begin{minted}[linenos,breaklines,frame=single]{latex}
\section{Prodotti notevoli}
\[
(a+b)^2=a^2+2ab+b^2
\]
\end{minted}
\end{document}
\end{lstlisting}
Il linguaggio indicato è \texttt{latex}.
L'opzione \texttt{linenos} numera le righe;
\texttt{breaklines} permette di spezzare quelle lunghe.

\prova
\begin{enumerate}
 \item Crea un progetto di prova e compila il modello.
 \item Sostituisci il contenuto con il codice di una tua equazione.
 \item Confronta il risultato con un blocco \texttt{lstlisting}.
\end{enumerate}
\controlla Se l'ambiente segnala un problema relativo all'elaboratore
esterno, conserva il messaggio e mostra il codice con \texttt{listings}.
Annota che cosa è riuscito a compilare nel tuo progetto.

\clearpage
\section{Indice analitico}\label{sec:analitico}
\index{indice!analitico}
L'indice analitico raccoglie parole e numeri di pagina.
L'indice generale raccoglie i titoli.
Alcune parole del laboratorio sono già marcate nel sorgente.
\begin{lstlisting}
% Nel preambolo:
\usepackage{makeidx}
\makeindex
% Nel corpo:
Un termine importante.\index{termine}
% Alla fine:
\printindex
\end{lstlisting}
La scrittura \verb|\index{indice!analitico}| crea una sottovoce.
Il passaggio di compilazione \texttt{makeindex} elabora il file
\texttt{.idx}; la successiva compilazione inserisce l'indice nel PDF.

\prova
\begin{enumerate}
 \item Trova nel sorgente le voci già inserite.
 \item Marca tre nuovi termini nei rispettivi paragrafi.
 \item Aggiungi due sottovoci sotto un termine comune.
 \item Aggiorna l'indice con il sistema di compilazione disponibile.
\end{enumerate}
\controlla Confronta ogni numero con la pagina in cui hai inserito
il comando. Se l'indice non appare, verifica il passaggio
\texttt{makeindex} nel registro di compilazione.

\clearpage
\section{Caratteri speciali, errori e avvisi}\label{sec:errori}
\subsection{Caratteri da scrivere esplicitamente}
\begin{center}
\begin{tabular}{cc}
\toprule
Risultato & Codice\\
\midrule
\% & \verb|\%|\\
\& & \verb|\&|\\
\_ & \verb|\_|\\
\# & \verb|\#|\\
\$ & \verb|\$|\\
\{ \} & \verb|\{ \}|\\
\textbackslash & \verb|\textbackslash|\\
\bottomrule
\end{tabular}
\end{center}
Questi caratteri svolgono funzioni particolari nel sorgente.
Per stamparli nel testo occorre usare una scrittura appropriata.

\subsection{Prove controllate}
\prova
\begin{enumerate}
 \item Scrivi una frase con una percentuale e un nome contenente un trattino basso.
 \item In una copia del progetto togli la graffa finale da
 \verb|\textbf{prova}|. Compila, leggi il primo errore e ripristina il comando.
 \item Cambia il nome di un richiamo senza cambiare l'etichetta.
 Osserva l'avviso e ripristina il riferimento.
 \item Allunga una parola in una cella finché supera il margine.
 Osserva l'avviso e rendi la colonna adatta al testo.
\end{enumerate}
\controlla Un PDF prodotto può ancora avere avvisi o riferimenti incompleti.
Correggi la causa e compila nuovamente.

\clearpage
\section{Progetto finale e controllo del lavoro}\label{sec:finale}
Prepara dentro la tua copia una breve relazione intitolata
\emph{Tre modi di studiare un triangolo}.
Scegli esempi dal laboratorio e adattali alla tua spiegazione.

\begin{enumerate}
 \item Organizza il testo in tre sottosezioni su lati, altezze e area.
 \item Inserisci due paragrafi, un elenco e una nota.
 \item Prepara una tabella con dati numerici di due triangoli.
 \item Disegna un triangolo, aggiungi didascalia e richiamalo nel testo.
 \item Inserisci una formula numerata per il calcolo dell'area e richiamala.
 \item Aggiungi un termine all'indice analitico.
\end{enumerate}

\subsection{Controllo prima della consegna}
\begin{itemize}
 \item Titolo, autore e indice sono aggiornati.
 \item Titoli ed elenchi hanno numerazione automatica.
 \item Le etichette hanno nomi unici e i richiami funzionano.
 \item Tabelle e figure stanno nei margini e hanno didascalie pertinenti.
 \item Le formule corrispondono ai dati delle figure.
 \item Hai ricompilato dopo l'ultima modifica.
\end{itemize}

\subsection{Una breve nota personale}
Scrivi quali esempi hai riutilizzato, che cosa hai cambiato
e un problema che hai risolto.
Se hai usato Gemini, indica una domanda precisa che gli hai posto
e una verifica che hai effettuato sul suggerimento.
La consegna comprende il sorgente e il PDF della tua versione.

\clearpage
\section*{Riferimenti per approfondire}
\phantomsection
\addcontentsline{toc}{section}{Riferimenti per approfondire}
Le spiegazioni, gli esempi matematici e le attività di questo laboratorio
sono stati preparati per la 2BS seguendo gli argomenti del centro
di apprendimento Overleaf.
Le guide originali permettono di approfondire i singoli comandi.

\begin{itemize}
 \item Indice degli argomenti\\
 \url{https://www.overleaf.com/learn/latex/Contents}
 \item Struttura\\
 \url{https://www.overleaf.com/learn/latex/Sections_and_chapters}
 \item Paragrafi e righe\\
 \url{https://www.overleaf.com/learn/latex/Paragraphs_and_new_lines}
 \item Elenchi\\
 \url{https://www.overleaf.com/learn/latex/Lists}
 \item Tabelle\\
 \url{https://www.overleaf.com/learn/latex/Tables}
 \item Posizione di figure e tabelle\\
 \url{https://www.overleaf.com/learn/latex/Positioning_images_and_tables}
 \item Intestazioni e piè di pagina\\
 \url{https://www.overleaf.com/learn/latex/Headers_and_footers}
 \item Contatori\\
 \url{https://www.overleaf.com/learn/latex/Counters}
 \item Indici analitici\\
 \url{https://www.overleaf.com/learn/latex/Indices}
 \item Evidenziazione del codice con minted\\
 \url{https://www.overleaf.com/learn/latex/Code_Highlighting_with_minted}
\end{itemize}

\clearpage
\phantomsection
\addcontentsline{toc}{section}{Indice analitico}
\printindex
\end{document}

